%!TEX root = ../template.tex

%% ----------------------------------------------------------------------------
%%  novathesis — 2-MainMatter/chapter-styles-showcase.tex
%% ----------------------------------------------------------------------------
%%
%% DEMO ONLY. Cycles through every available chapter style, one short
%% \chapter{} each, so they can be visually compared side by side. Not part
%% of the default build -- see the commented-out \ntaddfile line in
%% 0-Config/4_files.tex; enable it locally to eyeball chapter-style changes.
%%
%% Two clearly separate groups, in this order:
%%   1. NOVATHESIS-CUSTOM styles -- novathesisFiles/ChapStyles/*.sty, selected
%%      in a real thesis with \ntsetup{style/chapter=<name>}. \ntshowcase-
%%      chapterstyle below \input's each file directly, exactly what
%%      nt-render-matter.sty does for that key, just repeated once per style
%%      instead of once for the whole document.
%%   2. STANDARD MEMOIR styles -- built into memoir.cls itself, selected in a
%%      real thesis with a plain \chapterstyle{<name>} call (novathesis's own
%%      style/chapter key only knows its own ChapStyle files, not these).
%%      \ntshowcasememoirstyle below just calls \chapterstyle{<name>}
%%      directly -- no \input, no file-loading quirks to work around.
%%
%% Both groups reset memoir's chapterstyle state back to its stock ("default"
%% chapterstyle) values before every switch, via the internal \@chs@def@ult --
%% the same reset \makechapterstyle silently runs before applying any named
%% style's own clauses. Doing it explicitly here is redundant for any style
%% that goes through that machinery (all of them, including every custom one
%% below) but is kept as a documented safety net.
%%
%% CAVEAT: isel, fmv and compact (group 1) additionally tweak section-heading
%% fonts/numbering depth as a side effect (\setsecheadstyle, \setsecnumdepth,
%% ...). Those live outside memoir's chapterstyle system, so the reset above
%% does NOT undo them -- they will persist into any later \chapter/\section
%% in the same build. Keep this file last in the build order if that matters.
%% ----------------------------------------------------------------------------

\typeout{NT FILE 2-MainMatter/chapter-styles-showcase.tex}%

\makeatletter
% Reset every macro/length memoir's chapterstyle system controls back to
% memoir's own stock values -- the same internal reset \makechapterstyle
% silently runs before applying any named style's own clauses.
%
% Three things memoir's own 'culver' style (below) deliberately touches are
% NOT part of that reset, because they aren't scoped to a single chapter's
% appearance at all -- and since nothing about \chapterstyle ever restores
% them, culver's choices otherwise leak into EVERY later chapter for the
% rest of the document, not just culver's own:
%   \clearforchapter -- the macro \chapter itself calls to decide whether
%     to start a new page (normally \cleartorecto/\clearpage/\cleartoverso
%     depending on the openright/openany/openleft class option); culver
%     sets it to {} for a "no new page" continuous layout. This is what
%     made chapters stop starting on a new page after 'chappell': the very
%     next style in the list, 'culver', is the one that actually flips it.
%   \thechapter -- culver switches chapter numbering to Roman numerals.
%   \ps@chapter -- \aliaspagestyle{chapter}{headings} makes chapter-opening
%     pages show a running head; memoir's own default is 'plain' (none).
% Capture the document's real values once, up front, and restore them on
% every reset below, so culver's choices stay scoped to its own chapter.
\let\ntshowcase@clearforchapter\clearforchapter
\let\ntshowcase@thechapter\thechapter
\let\ntshowcase@pschapter\ps@chapter
\newcommand{\ntshowcaseresetchapterstyle}{%
  \@chs@def@ult
  \let\clearforchapter\ntshowcase@clearforchapter
  \let\thechapter\ntshowcase@thechapter
  \let\ps@chapter\ntshowcase@pschapter
}
\makeatother

% ChapStyle files are normally \input while the class is still loading,
% where '@' is already a letter; \input from the document body needs
% \makeatletter itself, since several styles use @-prefixed internals
% (\chs@ell@helper, \@nt@..., \@chapapp, ...). This must wrap the whole
% \newcommand definition below, not just its execution: '@'-containing
% tokens are tokenized once, when this file is read and the macro body is
% stored, so an inner \makeatletter (taking effect only when the macro
% later RUNS) would be too late to matter.
\makeatletter
\newcommand{\ntshowcasechapterstyle}[1]{%
  \ntshowcaseresetchapterstyle
  % The outer \makeatletter (above) only fixes how THIS file's own '@'
  % tokens were tokenized when this macro was DEFINED. \input below does
  % its OWN fresh tokenization of the ChapStyle file's characters, at
  % RUN time, so '@' must ALSO be made a letter again right here, live,
  % for that nested read -- an executable effect, not just a definition-
  % time one, hence still needed despite the outer wrap.
  \makeatletter
  % ist/ist2 (and elegant) do \RequirePackage{soul} -- harmless while the
  % class is loading (the normal case), but re-issuing it from the document
  % body prints a stray "soul" onto the page even though the package is
  % already preloaded for exactly this purpose (see 0-Config/5_packages.tex)
  % and the reload itself is a no-op. Silence it just for this \input --
  % NOT inside a \begingroup/\endgroup, since \chapterstyle's own
  % \renewcommand's must survive past this macro, not just this group.
  \let\ntshowcase@RequirePackage\RequirePackage
  \let\RequirePackage\@gobble
  \input{\DIRCHAPSTYLES/#1.sty}%
  \let\RequirePackage\ntshowcase@RequirePackage
  \makeatother
  \chapter{Example Title in Chapter style: \texttt{#1}}
  \label{cha:style-showcase-#1}
  This chapter demonstrates the \texttt{#1} chapter style, selected with:

  \texttt{\textbackslash ntsetup\{style/chapter=#1\}}

  \lipsum[1]
}
\makeatother

%%%%%%%%%%%%%%%%%%%%%%%%%%%%%%%%%%%%%%%%%%%%%%%%%%%%%%%%%%%%%%%%%%%%%%%%
%% 1. NOVATHESIS-CUSTOM CHAPTER STYLES
%%    \ntsetup{style/chapter=<name>}
%%%%%%%%%%%%%%%%%%%%%%%%%%%%%%%%%%%%%%%%%%%%%%%%%%%%%%%%%%%%%%%%%%%%%%%%
\ntshowcasechapterstyle{bar}
\ntshowcasechapterstyle{bar-compact}
\ntshowcasechapterstyle{bluebox}
\ntshowcasechapterstyle{compact}
% 'elegant' is skipped here: it has a pre-existing bug independent of this
% showcase -- \so{\@chapapp} (soul's letter-spacing wrapped around the
% chapter/appendix/annex name) corrupts babel's \bbl@ensure@english state
% and crashes the *next* \chapter call, reproducible even as the document's
% only configured style/chapter (i.e. not a showcase artifact). See the
% separately-flagged bug report.
% \ntshowcasechapterstyle{elegant}
\ntshowcasechapterstyle{fmv}
\ntshowcasechapterstyle{isel}
\ntshowcasechapterstyle{ist}
\ntshowcasechapterstyle{ist2}
\ntshowcasechapterstyle{vz34}
\ntshowcasechapterstyle{vz43}


%%%%%%%%%%%%%%%%%%%%%%%%%%%%%%%%%%%%%%%%%%%%%%%%%%%%%%%%%%%%%%%%%%%%%%%%
%% 2. STANDARD MEMOIR CHAPTER STYLES
%%    \chapterstyle{<name>}  (novathesis does not wrap these -- there is no
%%    style/chapter value for them; call \chapterstyle{<name>} yourself,
%%    e.g. right after \documentclass or in your own \ntaddfile'd chapter)
%%%%%%%%%%%%%%%%%%%%%%%%%%%%%%%%%%%%%%%%%%%%%%%%%%%%%%%%%%%%%%%%%%%%%%%%
\newcommand{\ntshowcasememoirstyle}[1]{%
  \ntshowcaseresetchapterstyle
  \chapterstyle{#1}
  \chapter{Example Title in Memoir style: \texttt{#1}}
  \label{cha:style-showcase-memoir-#1}
  This chapter demonstrates memoir's built-in \texttt{#1} chapter style,
  selected with:

  \texttt{\textbackslash chapterstyle\{#1\}}

  \lipsum[1]
}

% Every \makechapterstyle{<name>}{...} memoir.cls itself defines, in the
% order they appear there. 'default' is memoir's own bare-bones fallback
% (no visual flourish -- it exists as the base every other style, including
% novathesis's own, resets to before applying its own overrides); 'section'
% and 'ell' are the two memoir styles novathesis's own 'compact'/'fmv' and
% 'bar'/'bar-compact' above are themselves built on -- shown here unpatched,
% for comparison against those.
\ntshowcasememoirstyle{default}
\ntshowcasememoirstyle{section}
\ntshowcasememoirstyle{article}
\ntshowcasememoirstyle{reparticle}
\ntshowcasememoirstyle{hangnum}
\ntshowcasememoirstyle{companion}
\ntshowcasememoirstyle{demo}
\ntshowcasememoirstyle{bianchi}
\ntshowcasememoirstyle{bringhurst}
\ntshowcasememoirstyle{brotherton}
\ntshowcasememoirstyle{chappell}
\ntshowcasememoirstyle{culver}
\ntshowcasememoirstyle{dash}
\ntshowcasememoirstyle{demo2}
\ntshowcasememoirstyle{demo3}
\ntshowcasememoirstyle{ell}
\ntshowcasememoirstyle{ger}
\ntshowcasememoirstyle{lyhne}
\ntshowcasememoirstyle{madsen}
\ntshowcasememoirstyle{pedersen}
\ntshowcasememoirstyle{southall}
\ntshowcasememoirstyle{thatcher}
\ntshowcasememoirstyle{veelo}
\ntshowcasememoirstyle{verville}
\ntshowcasememoirstyle{crosshead}
\ntshowcasememoirstyle{dowding}
\ntshowcasememoirstyle{komalike}
\ntshowcasememoirstyle{ntglike}
\ntshowcasememoirstyle{tandh}
\ntshowcasememoirstyle{wilsondob}


% Restore whatever chapter style the document is actually configured to use,
% so any real content coming after this demo chapter is not left showing
% the last showcased style instead of the user's chosen one.
\ntshowcaseresetchapterstyle
\makeatletter
\InputIfFileExists{\DIRCHAPSTYLES/\option{/novathesis/style/chapter}.sty}
  {}{\chapterstyle{\option{/novathesis/style/chapter}}}
\makeatother
